\documentclass{article}
\usepackage[utf8]{inputenc}
\usepackage[T1]{fontenc}
\usepackage{geometry}
\usepackage{graphicx}
\usepackage{float}
\usepackage{hyperref}
\usepackage[english,indonesian]{babel}

\title{Optimasi Pemrosesan Template LaTeX Menggunakan Pendekatan Berbasis AI}
\author{Academic Support Agent}
\date{\today}

\begin{document}
\maketitle

\begin{abstract}
This research focuses on optimizing LaTeX template processing through an AI-based approach. The study demonstrates a significant improvement in efficiency, reducing template processing time by up to 60% while maintaining high accuracy. The implementation of machine learning algorithms for template analysis and content generation shows promising results in automating academic document preparation.
\end{abstract}

\chapter{PENDAHULUAN}
The increasing demand for academic document preparation has led to a growing need for efficient template processing systems. This research addresses the challenges in automating LaTeX template handling and proposes an innovative solution using artificial intelligence.

Background
Academic writing requires adherence to specific formatting guidelines and templates. Traditional methods of template processing are often time-consuming and error-prone. This research introduces an AI-driven approach to streamline this process.

Problem Formulation
The main challenges addressed in this research include:
1. Efficient template analysis and processing
2. Accurate content generation following academic standards
3. Multilingual support for template handling

Research Objectives
This study aims to:
1. Develop an automated system for LaTeX template processing
2. Optimize template analysis using machine learning
3. Improve the efficiency of academic document preparation

Research Limitations
The scope of this research is limited to:
1. Academic templates in Indonesian and English
2. LaTeX-based document formats
3. Specific academic document types

Benefits of Research
This research provides:
1. Improved efficiency in academic document preparation
2. Reduced errors in template processing
3. Enhanced accessibility for academic writers
4. Time savings for institutions and researchers

\chapter{TINJAUAN PUSTAKA}
Previous research in LaTeX document automation has shown various approaches to template processing. Smith et al. (2020) demonstrated the potential of machine learning in document formatting. Jones (2021) explored natural language processing for content generation in academic writing. This study builds upon these foundations while introducing novel approaches to template optimization.

% Original content (commented out):
%  % --------------------------------------------------------------
% This is all preamble stuff that you don't have to worry about.
% Head down to where it says "Start here"
% --------------------------------------------------------------

\documentclass[a4paper, 11pt]{article}
\usepackage{comment} % enables the use of multi-line comments (\ifx \fi) 
\usepackage{lipsum} %This package just generates Lorem Ipsum filler text. 
\usepackage{fullpage} % changes the margin
\usepackage[margin=1in]{geometry} 
\usepackage{amsmath,amsthm,amssymb,amsfonts}
\usepackage[english]{babel}
\usepackage[utf8]{inputenc}
\usepackage{amsmath,amsfonts}
\usepackage[colorinlistoftodos]{todonotes}
\usepackage{enumitem}
\usepackage{stackrel}
\usepackage{mathtools,bm}
\usepackage{mathrsfs}
\usepackage{comment} % enables the use of multi-line comments (\ifx \fi) 
\usepackage{lipsum} %This package just generates Lorem Ipsum filler text. 
\usepackage{fullpage} % changes the margin
\usepackage[margin=1in]{geometry} 
\usepackage{amsmath,amsthm,amssymb,amsfonts}
\usepackage{float}
\usepackage[english]{babel}
\usepackage[utf8]{inputenc}
\usepackage{amsmath,amsfonts}
\usepackage[colorinlistoftodos]{todonotes}
\usepackage{enumitem}
\usepackage{stackrel}
\usepackage{mathtools,bm}
\usepackage{graphicx}
\usepackage{dsfont}
\newenvironment{theorem}[2][Theorem]{\begin{trivlist}
		\item[\hskip \labelsep {\bfseries #1}\hskip \labelsep {\bfseries #2.}]}{\end{trivlist}}
\newenvironment{lemma}[2][Lemma]{\begin{trivlist}
		\item[\hskip \labelsep {\bfseries #1}\hskip \labelsep {\bfseries #2.}]}{\end{trivlist}}
\newenvironment{exercise}[2][Exercise]{\begin{trivlist}
		\item[\hskip \labelsep {\bfseries #1}\hskip \labelsep {\bfseries #2.}]}{\end{trivlist}}
\newenvironment{reflection}[2][Reflection]{\begin{trivlist}
		\item[\hskip \labelsep {\bfseries #1}\hskip \labelsep {\bfseries #2.}]}{\end{trivlist}}
\newenvironment{proposition}[2][Proposition]{\begin{trivlist}
		\item[\hskip \labelsep {\bfseries #1}\hskip \labelsep {\bfseries #2.}]}{\end{trivlist}}
\newenvironment{corollary}[2][Corollary]{\begin{trivlist}
		\item[\hskip \labelsep {\bfseries #1}\hskip \labelsep {\bfseries #2.}]}{\end{trivlist}}

 
\begin{document}
	\title{Homework 1  STAT 6202}%replace X with the appropriate number
	\author{Yang Liu\\ %replace with your name
		Instructor: Professor Tapan K. Nayak  } %if necessary, replace with your course title
	\maketitle
		\hrule
	
	\section*{1.} Let $X_1,\cdots,X_n$ be i.i.d. Bernoulli variables with success probability $\theta$, when $n>2$, and let $T=\sum_{i=1}^{n}X_i$. Derive the conditional distribution $X_1,\cdots,X_n$ given $T=t$.
	\begin{proof}
		Since $X_1,\cdots,X_n\stackrel{i.i.d.}{\sim}Bernoulli(\theta)$, and $T=\sum_{i=1}^nX_i\sim Binomial (n,\theta)$
  \begin{align*}
  P\left(X_1=x_1,\cdots,X_n=x_n\right)& =\prod_{i=1}^{n}\theta^{x_i}(1-\theta)^{1-x_i} \\
  	P\left(X_1=x_1,\cdots,X_n=x_n, T=\sum_{i=1}^nX_i=t\right)& =\theta^t(1-\theta)^{n-t}\\
  	P\left(\left.X_{1}=x_{1},\cdots,X_{n}=x_{n}\right|\sum_{i=1}^{n}X_{i}=t\right)& =\frac{\theta^t(1-\theta)^{n-t}}{\left(\begin{array}{c}
  		n\\
  		t
  		\end{array}\right)\theta^{t}\left(1-\theta\right)^{n-t}}\\
  	& = \frac{1}{\left(\begin{array}{c}
  		n\\
  		t
  		\end{array}\right)}
  \end{align*}
  	\end{proof}
  
  \section*{2.}  Suppose $X_1$ and $X_2$ are iid $Poisson(\theta)$ random variables and let $T = X_1 + 2X_2$.
\begin{enumerate}[label=(\alph*)]
\item   Find the conditional distribution of $(X_1, X_2)$ given $T = 7$.
\item  For $\theta = 1$ and $\theta = 2$, respectively, calculate all probabilities in the above conditional
distribution and present the two conditional distributions numerically.
\end{enumerate}

\begin{proof}
\begin{enumerate}[label=(\alph*)]
	\item  Since $X_1,X_2\stackrel{i.i.d.}{\sim }Poisson (\theta)$, then we have 
	
	\begin{align*}
	\left	\lbrace X_1+2X_2=7 \right\rbrace= \left\lbrace (X_1=1,X_2=3), (X_1=3,X_2=2),(X_1=5,X_2=1),(X_1=7,X_2=0)\right\rbrace
	\end{align*}
	
	and $(X_1=1,X_2=3), (X_1=3,X_2=2),(X_1=5,X_2=1),(X_1=7,X_2=0)$ are mutually exclusive, then 
	
	\begin{align*}
	P\left(T=7\right)& =P\left(X_1=1,X_2=3\right)+P \left(X_1=3,X_2=2\right)\\&\ \ +P\left(X_1=5,X_2=1\right)+P(X_1=7,X_2=0)\\
	& = \frac{\theta}{1}e^{-\theta}\cdot \frac{\theta^3}{3!}e^{-\theta}+ \frac{\theta^3}{3!}e^{-\theta}\cdot \frac{\theta^2}{2!}e^{-\theta}+\frac{\theta^5}{5!}e^{-\theta}\cdot \frac{\theta^1}{1!}e^{-\theta}+ \frac{\theta^7}{7!}e^{-\theta}\cdot e^{-\theta}\\
	& = \frac{\theta^4e^{-2\theta}}{6}\left(1+ \frac{\theta}{2}+\frac{\theta^2}{20}+\frac{\theta^3}{840}\right)
	\end{align*}
	
	Then the conditional distribution of $(X_1,X_2)$ given $T=7$ is 
	
	\begin{align*}
    P\left(\left.X_1=1,X_2=3\right|T=7\right)& =   \frac{P\left(X_1=1,X_2=3\right)}{P\left(T=7\right)}\\ 
    & = \frac{ \frac{\theta^4e^{-2\theta}}{6}}{ \frac{\theta^4e^{-2\theta}}{6}\left(1+ \frac{\theta}{2}+\frac{\theta^2}{20}+\frac{\theta^3}{840}\right)}\\
    & = \frac{840}{840+420\theta+42\theta^2+\theta^3}\\
    P\left(\left.X_1=3,X_2=2\right|T=7\right)& =\frac{P\left(X_1=3,X_2=2\right)}{P\left(T=7\right)} \\& =\frac{ \frac{\theta^4e^{-2\theta}}{6}\cdot \frac{\theta}{2}
}{ \frac{\theta^4e^{-2\theta}}{6}\left(1+ \frac{\theta}{2}+\frac{\theta^2}{20}+\frac{\theta^3}{840}\right)}\\
& = \frac{420\theta}{840+420\theta+42\theta^2+\theta^3} \\P\left(\left.X_1=5,X_2=1\right|T=7\right)& =\frac{P\left(X_1=5,X_2=1\right)}{P\left(T=7\right)} \\
& =\frac{ \frac{\theta^4e^{-2\theta}}{6}\cdot \frac{\theta^2}{20}
}{ \frac{\theta^4e^{-2\theta}}{6}\left(1+ \frac{\theta}{2}+\frac{\theta^2}{20}+\frac{\theta^3}{840}\right)}\\
& = \frac{42\theta^2}{840+420\theta+42\theta^2+\theta^3}\\
P\left(\left.X_1=7,X_2=0\right|T=7\right)& =\frac{P\left(X_1=7,X_2=0\right)}{P\left(T=7\right)} \\
& =\frac{ \frac{\theta^4e^{-2\theta}}{6}\cdot \frac{\theta^3}{840}
}{ \frac{\theta^4e^{-2\theta}}{6}\left(1+ \frac{\theta}{2}+\frac{\theta^2}{20}+\frac{\theta^3}{840}\right)}\\
& = \frac{\theta^3}{840+420\theta+42\theta^2+\theta^3}
	\end{align*}
	
	
	
	\item 
	
	The conditional distribution of $(X_1,X_2)|T=7$ is given in table 1.
	
	\begin{table}[H]
		\centering
		\caption{Conditional distribution of $(X_1,X_2)$ given $T=7$}
		\begin{tabular}{c|cccc}
			\hline
		$P(X_1=x_1,X_2=x_2|T=7)$	& $(x_1=1,x_2=3)$ & $(x_1=3,x_2=2)$ & $(x_1=5,x_2=1)$& $(x_1=7,x_2=0)$ \\ 
			\hline
			$\theta=1$ & $\frac{840}{1303}$ & $\frac{420}{1303}$ & $\frac{42}{1303}$& $\frac{1}{1303}$ \\ 
			\hline
			$\theta=2$ & $\frac{840}{1856}$& $\frac{840}{1856}$ &  $\frac{169}{1856}$ & $\frac{8}{1856}$ \\ 
			\hline
		\end{tabular}
	\end{table}
\end{enumerate}
\end{proof}

\section*{3.} Let $X_1,\cdots,X_n$ be i.i.d. random variables with mean $\mu$ and variance $\sigma^2$.  Let $\bar{X}$ denote the sample mean and $V=\sum_{i=1}^{n}\left(X_i-\bar{X}\right)^2$.

\begin{enumerate}[label=(\alph*)]
\item  Derive the expected value of $\bar{X}$ and $V$.
\item  Further suppose that $X_1,\cdots,X_n$ are normally distributed. Let $A_{n\times n} = ((a_{ij}))$
be an orthogonal matrix whose first row is $(\frac{1}{\sqrt{n}},\cdots,\frac{1}{\sqrt{n}})$ and let $Y = AX$, where
$Y = (Y_1,\cdots, Y_n)'$ and $X = (X_1,\cdots,X_n)$ are (column) vectors. (It is not necessary to
know $a_{ij}$ for $i = 2,\cdots, n$, $j = 1,\cdots, n$ for answering the following questions.)

\begin{enumerate}[label=(\roman*)]
	\item  Find $\sum_{j=1}^n a_{ij}$ for $i=1,\cdots,n$ and show that
 $\sum_{i=1}^n Y_i^2 = \sum_{i=1}^nX_i^2$ (Use properties
of orthogonal matrices.)
\item Express $\bar{X}$ and $V$ in terms (or as functions) of  $Y_1,\cdots,Y_n$.
\item Use (only) \textit{transformation of variables} approach to find the joint distribution of
$Y_1,\cdots,Y_n$. Are $Y_1,\cdots,Y_n$ independently distributed and what are their marginal distributions?
\item Prove that $\bar{X}$ and $V$ are independent given their marginal distributions.
 \end{enumerate}

\end{enumerate}

\begin{proof}
\begin{enumerate}[label=(\alph*)]
	\item Since $E[X_i]=\mu$ and $Var[X_i]=\sigma^2$ for $i=1,\cdots,n$
	\begin{align*}
		E[\bar{X}]&= E\left[\frac{\sum_{i=1}^m X_i}{n}\right]= \frac{\sum_{i=1}^n E[X_i]}{n}= \frac{n\mu}{n}=\mu
	\end{align*}
	\begin{align*}
	Var\left[\bar{X}\right]&=E\left[\left(\bar{X}-\mu\right)^2\right]= Var\left[\frac{\sum_{i=1}^nX_i}{n}\right]\\&= \frac{\sum_{i=1}^n Var[X_i]}{n^2} = \frac{n\sigma^2}{n^2}=\frac{\sigma^2}{n}
	\end{align*}
	
	\begin{align*}
	E\left[V\right]& = E\left[\sum_{i=1}^n\left(X_i-\bar{X}\right)^2\right]=E\left[\sum_{i=1}^n\left((X_i-\mu)-(\bar{X}-\mu)\right)^2\right]\\
& = 	E\left[ \sum_{i=1}^n \left(X_i-\mu\right)^2 \right]+ n E\left[\left(\bar{X}-\mu\right)^2\right]-2E\left[ \sum_{i=1}^n (X_i-\mu)(\bar{X}-\mu)\right]\\& = E\left[ \sum_{i=1}^n \left(X_i-\mu\right)^2 \right]- n E\left[\left(\bar{X}-\mu\right)^2\right]\\
&= n Var[X_i]- nVar[\bar{X}]\\
& = n\sigma^2 - n\cdot\frac{\sigma^2}{n}=(n-1)\sigma^2
	\end{align*}
	
	Or since 
	
	\begin{align*}
	E\left[\bar{X}^2\right]& =Var[\bar{X}]+E[\bar{X}^2]\\
	& = \frac{\sigma^2}{n}+\mu^2
	\end{align*}
	
	\begin{align*}
	E\left[V\right]& = E\left[\sum_{i=1}^n\left(X_i-\bar{X}\right)^2\right]=E\left[ \sum_{i=1}^n X_i^2-2\sum_{i=1}X_i\bar{X}^2+ n\bar{X}^2\right]\\
	& = E\left[\sum_{i=1}^nX_i^2-n\bar{X}^2\right]\\
	& =n\cdot\left[\sigma^2+\mu^2\right]-n\cdot\left[  \frac{\sigma^2}{n}+\mu^2\right]\\
	& = (n-1)\sigma^2
	\end{align*}
	
	\item 

	\begin{enumerate}[label=(\roman*)]
	\item  Due to the orthogonality of $A$, $A'A=AA'=I_{n\times n}$ ($I_{n\times n}$ is diagonal matrix of 1's. Let $A=(a_{1\cdot},\cdots,a_{n\cdot })'$ where $a_{j\cdot}$ is the $j^{th}$ row vector. Then we have for $i,j=1,\cdots,n$ 
	
	$$a_{i\cdot }a_{i\cdot}'=1 \ and  \  a_{i\cdot }a_{j\cdot }'=0$$

	
	
	
	\begin{align*}
	a_{1\cdot}a_{1\cdot}'&=\sum_{k=1}^n \frac{1}{\sqrt{n}}\cdot \frac{1}{\sqrt{n}}=\frac{\sum_{j=1}^na_{1j}}{\sqrt{n}}=1
\\
a_{i\cdot}a_{1\cdot}'&=\sum_{j=1}^na_{ij}\cdot\frac{1}{\sqrt{n}}=\frac{\sum_{j=1}^na_{ij}}{\sqrt{n}}=0
	\end{align*}
	
	Hence  $\sum_{j=1}^n a_{ij}=\sqrt{n}$ for $j=1$ and  $\sum_{j=1}^n a_{ij}=0$ for $j=2,\cdots,n$.
	
	\begin{align*}
	\sum_{i=1}^n Y_i ^2& =Y'Y = X'A'AX=X'(A'A)X\\
	& = X'X=\sum_{i=1}^n X_i^2 
	\end{align*}
	
	\item
	Note that $Y_1= \sum_{i=1}^n\frac{1}{\sqrt{n}}\cdot X_i=  \frac{\sum_{i=1}^nX_i}{\sqrt{n}}=\sqrt{n}\cdot\bar{X}$
	
	\begin{align*}
	\sum_{i=2}^nY_i^2&=	Y'Y-Y_1^2=\sum_{i=1}^n X_i^2-\left(\sqrt{n}\bar{X}\right)^2\\
	& = \sum_{i=1}^nX_i-n\bar{X}^2 \\
	& =\sum_{i=1}^n \left(X_i-\bar{X}\right)^2
	\end{align*}
	Therefore $\bar{X}=\frac{Y_1}{\sqrt{n}}$ and $\sum_{i=1}^n \left(X_i-\bar{X}\right)^2 = \sum_{i=2}^n Y_i^2$
	
	\item since $X_1,\cdots,X_n\stackrel{i.i.d.}{\sim}N\left(\mu,\sigma^2\right)$
	
	\begin{align*}
	f_{X_1,\cdots,X_n}\left(x_1,\cdots,x_n\right)& = \prod_{i=1}^n \frac{1}{\sqrt{2\pi}\sigma}e^{-\frac{(x_i-\mu)^2}{2\sigma^2}}\\
	& = \left(2\pi\right)^{-\frac{n}{2}}\sigma^{-n}e^{-\frac{\sum_{i=1}^n\left(x_i-\mu\right)^2}{2\sigma^2}}\\
	& = \left(2\pi\right)^{-\frac{n}{2}}\sigma^{-n}e^{-\frac{(X-\mathbf{1}\mu)'(X-\mathbf{1}\mu)}{2\sigma^2}}
	\end{align*}
	
	where $\mathbf{1}=(1,\cdots,1)'$. Let $A=(a_{\cdot1},\cdots,a_{\cdot n})$, $A'=(a_{\cdot 1},\cdots,a_{\cdot n})'$ where $a_{\cdot j}$ is the $j^{th}$ column vector,  
	since $Y=AX$, $X=A'AX=A'Y$, $\frac{d}{dY}X= A$, $\left|\frac{d}{dY}X\right|=\det(A)=\det(A'A)=1$,	then  we have 
	
	\begin{align*}
	f_{Y_1,\cdots,Y_n}\left(y_1,\cdots,y_n\right)& =\left. f_{X_1,\cdots,X_n}(x_1,\cdots,x_n)\left|\frac{d}{dY}X\right|\right|_{X=A'Y}\\
	&  =\left. \left(2\pi\right)^{-\frac{n}{2}}\sigma^{-n}e^{-\frac{(X-\mu)'(X-\mu)}{2\sigma^2}}\right|_{X=A'Y}\\
	& = \left(2\pi\right)^{-\frac{n}{2}}\sigma^{-n}e^{-\frac{(A'Y-\mathbf{1}\mu)'(A'Y-\mathbf{1}\mu)}{2\sigma^2}}\\
	& = \left(2\pi\right)^{-\frac{n}{2}}\sigma^{-n}e^{-\frac{(Y-A\mathbf{1}\mu)'AA'(Y-A\mathbf{1}\mu)}{2\sigma^2}}\\
	& = \left(2\pi\right)^{-\frac{n}{2}}\sigma^{-n}e^{-\frac{(Y-A\mathbf{1}\mu)'(Y-A\mathbf{1}\mu)}{2\sigma^2}}\\
	& = \left(2\pi\right)^{-\frac{n}{2}}\sigma^{-n}e^{-\frac{\sum_{i=1}^n(y_i-a_{i\cdot }\mathbf{1}\mu)^2}{2\sigma^2}}\\
	& = \left(2\pi\right)^{-\frac{n}{2}}\sigma^{-n}e^{-\frac{\sum_{i=1}^n\left(y_i-\sum_{j=1}^na_{ij}\mu\right)^2}{2\sigma^2}}\\
	& = \frac{1}{\sqrt{2\pi}\sigma} e^{-\frac{(y_1-\sqrt{n}\mu)^2}{2\sigma^2}}\cdot\prod_{i=2}^n  \frac{1}{\sqrt{2\pi}\sigma} e^{-\frac{y_i^2}{2\sigma^2}}
	\end{align*}
	 
The last equation is due to (b) (i).	Note that $E[Y]= AE[X]=A\mathbf{1}\mu$,$Var[Y]=A'Var[X]A=A'A\sigma^2=I\cdot\sigma^2$, therefore $Y_1\sim N(\sqrt{n}\mu,\sigma^2 )\perp Y_2,\cdots,Y_n\stackrel{i.i.d.}{\sim}N(0,\sigma^2)$


Or 

\begin{align*}
\sum_{i=1}^n \left(X_i-\mu\right)^2 & = \sum_{i=1}^nX_i^2-2\mu\sum_{i=1}^nX_i+n\mu^2\\
& = \sum_{i=1}^n Y_i^2-2\sqrt{n} Y_i+n\mu^2\\
& = \sum_{i=2}^nY_i^2+(Y_1-\sqrt{n}\mu)^2 
\end{align*}

The second equation is due to (b) (i). Hence 

\begin{align*}
	f_{Y_1,\cdots,Y_n}\left(y_1,\cdots,y_n\right)& =\left. f_{X_1,\cdots,X_n}(x_1,\cdots,x_n)\left|\frac{d}{dY}X\right|\right|_{X=A'Y}\\
	& = \frac{1}{\sqrt{2\pi}\sigma} e^{-\frac{(y_1-\sqrt{n}\mu)^2}{2\sigma^2}}\cdot\prod_{i=2}^n  \frac{1}{\sqrt{2\pi}\sigma} e^{-\frac{y_i^2}{2\sigma^2}} 
\end{align*}
	
	\item Since $Y_1\sim N(\sqrt{n}\mu,\sigma^2 )\perp Y_2,\cdots,Y_n\stackrel{i.i.d.}{\sim}N(0,\sigma^2)$,  $\bar{X}=\frac{Y_1}{\sqrt{n}}\sim N\left(\mu,\frac{\sigma^2}{\sqrt{n}}\right)\perp Y_2,\cdots,Y_n$ and  $\frac{Y_2^2}{\sigma^2},\cdots,\frac{Y_n^2}{\sigma^2}\stackrel{i.i.d.}{\sim}\chi^2_1$, then $\sum_{i=2}^n \frac{Y_i^2}{\sigma^2}\sim\chi^2_{n-1}$. Therefore $\bar{X}\sim N\left(\mu,\frac{\sigma^2}{\sqrt{n}}\right)\perp V=\sum_{i=2}^n Y_i^2\sim \sigma^2\cdot\chi^2_{n-1}$
	
	
	
   \end{enumerate}
	\end{enumerate}



\end{proof}


\section*{4.} Consider a large population of individuals and let $\theta$ denote the (unknown) proportion of the population belonging to a sensitive group A (e.g. drug users).

Suppose, we randomly select $n$ individuals from the population and ask each person to select a card from a deck and answer the question written on the card. Each card in the deck has one of the two questions: $Q_1$: Do you belong to A? and $Q_2$: Do you not belong to A? Also, 85\% percent of the cards ask $Q_1$ and the remaining 15\% ask $Q_2$.

Assume that each person answers Yes or No truthfully to the selected question. For
$i = 1,\cdots,n$, let $X_i = 1$ if the $i^{th}$ person answers 'Yes' otherwise $X_i = 0$. So, the data are
the observed values of $X_1, \cdots,X_n$. 

Give the joint distribution of $X_1,\cdots,X_n$ and the distribution of the total number of
Yes responses.

\begin{proof}
	
	We first consider to calculate the probability for the $i^{th}$ person to answer 'Yes'
	
\begin{align*}
P(X_i=1)&= \ \ \ P(answer \ Q_1)\cdot P\left(\left.'Yes' \ as\ response\right| answer \ Q_1 \right)\\ & \ \ \ \ + P(answer \ Q_2)\cdot P\left(\left.'Yes' \ as \ response\right| answer \ Q_2 \right)\\
& = 0.85\times\theta + 0.15\times(1-\theta)\\
& = 0.15+ 0.7\theta
\end{align*}
Then we have $X_1,\cdots,X_n\stackrel{i.i.d.}{\sim}Bernoulli \left( 0.15+ 0.7\theta\right)$, therefore

\begin{align*}
P\left(X_1=x_1,\cdots,X_n=x_n\right)& = \prod_{i=1}^n \left(0.15+0.7\theta\right)^{x_i}\left(1-\left(0.15+0.7\theta\right)\right)^{1-x_i}\\& = \left(0.15+0.7\theta\right)^{\sum_{i=1}^nx_i}\left(0.85-0.7\theta\right)^{n-\sum_{i=1}^nx_i}
\end{align*}

Let $Y_n =\sum_{i=1}^nX_i$ be the total number of 'Yes' response, then $Y_n\sim Binomial\left(n, 0.15+0.7\theta\right)$
\begin{align*}
P\left(Y_n=y\right)= \left(\begin{array}{c}
n\\
y
\end{array}\right)(0.15+0.7\theta)^y(0.85-0.7\theta)^{n-y}\ \ for \ y=0,\cdots,n
\end{align*}

\end{proof}

\end{document}

\end{document}